% !TEX root = task-1-10-4.tex
% Editable task bank. Append new BankTask blocks at the end of the document.
\documentclass[12pt,a4paper]{article}
\usepackage{physics-textbook}
\usepackage{graphicx,needspace}
\usetikzlibrary{decorations.pathmorphing}
\newif\ifShowSolutions
\ifdefined\TutorEdition
  \ShowSolutionstrue
\else
  \ShowSolutionsfalse % Change to \ShowSolutionstrue to display solutions.
\fi
\newcounter{banktask}
\newsavebox{\BankTaskBox}
\newcommand{\BankTaskContent}[3]{%
  \noindent\textbf{\thebanktask.}\quad #2\par
  \ifShowSolutions
    \Needspace{50mm}% Keep the solution heading with its opening text or diagram.
    #3\par
    \ifdefined\TutorEdition
      \ifcsname TutorSource#1\endcsname
        \nopagebreak\smallskip{\small\csname TutorSource#1\endcsname\par}%
      \fi
    \fi
  \fi
}
\newcommand{\BankTask}[3]{%
  \refstepcounter{banktask}\par\medskip
  \sbox{\BankTaskBox}{\begin{minipage}{\linewidth}%
    \BankTaskContent{#1}{#2}{#3}\end{minipage}}%
  \ifdim\dimexpr\ht\BankTaskBox+\dp\BankTaskBox\relax>\dimexpr\textheight-20mm\relax
    \Needspace{35mm}\BankTaskContent{#1}{#2}{#3}%
  \else
    \Needspace{\dimexpr\ht\BankTaskBox+\dp\BankTaskBox\relax}%
    \noindent\usebox{\BankTaskBox}\par
  \fi
  \medskip
}
\begin{document}
\BookSubsection{1.10.4}{Искусственные спутники}


% Source sdamgia-8408: https://phys-ege.sdamgia.ru/problem?id=8408
\BankTask{sdamgia-8408}{
Два искусственных спутника Земли массой $m_1 = 200\,\text{кг}$ и $m_2= 400\,\text{кг}$ обращаются по круговым орбитам одинакового радиуса. Чему равно отношение скоростей этих спутников $\frac{v _2}{v _1}$?
}{
\smallskip\noindent\textbf{Решение}\par
Пусть $M$ — масса Земли, $R$ — радиус орбиты, $v$ — скорость движения тела, $m$ — масса тела. При вращении тела по окружности вокруг Земли его центростремительное ускорение вызывается гравитационным притяжением Земли:

\[ma_{\text{ц}}=G\frac{mM}{R^2}\Longleftrightarrow m\frac{v^2}{R}=G\frac{mM}{R^2}\Longleftrightarrow\frac{v^2}{R}=G\frac{M}{R^2}\Longleftrightarrow v=\sqrt{\frac{GM}{R}}.\]

Заметим, что скорость не зависит от массы тела, а зависит только от радиуса вращения. Следовательно, скорости спутников, вращающихся по круговым орбитам одинакового радиуса равны: $\frac{v _2}{v _1} =1.$

Ответ: $1$.
}

% Source sdamgia-12851: https://phys-ege.sdamgia.ru/problem?id=12851
\BankTask{sdamgia-12851}{
Два искусственных спутника движутся вокруг однородной сферической планеты по круговым орбитам. Радиус орбиты первого спутника $800\,\text{км}$, масса этого спутника $50\,\text{кг}$. Радиус орбиты второго спутника $1600\,\text{км}$. При этом спутники притягиваются к планете с одинаковыми по модулю силами. Чему равна масса второго спутника? (В килограммах.)
}{
\smallskip\noindent\textbf{Решение}\par
Приравняем силы притяжения спутников к планете и выразим массу второго спутника:

\[\frac{Gm_1M}{R_1 ^2} = \frac{Gm_2M}{R_2 ^2} \Longleftrightarrow  \frac{m_1}{R_1 ^2} = \frac{m_2}{R_2 ^2} \Longleftrightarrow \]

\[\Longleftrightarrow  m_2=m_1 \cdot  \frac{R_2 ^2}{R_1 ^2} =50 \cdot  \frac{1600 ^2}{800 ^2} =200\,\text{кг}.\]

Ответ: $200$.
}

% Source sdamgia-32058: https://phys-ege.sdamgia.ru/problem?id=32058
\BankTask{sdamgia-32058}{
Ускорение свободного падения на поверхности Юпитера в $2{,}6$ раза больше, чем на поверхности Земли. Первая космическая скорость для Юпитера в $5{,}4$ раза больше, чем для Земли. Во сколько раз радиус Юпитера больше радиуса Земли? Ответ округлите до целого числа.
}{
\smallskip\noindent\textbf{Решение}\par
Первая космическая скорость на поверхности планеты равна $v=\sqrt{gR},$ откуда радиус планеты равен $R= \frac{v ^2}{g} .$

Таким образом, отношение радиусов Юпитера и Земли

\[\frac{R_{\text{ю}}}{R_{\text{з}}} = \frac{v _{\text{ю}} ^2 g_{\text{з}}}{v _{\text{з}} ^2 g_{\text{ю}}} = \frac{5{,}4 ^2}{2{,}6} \approx11.\]

Ответ: $11$.
}

% Source sdamgia-9306: https://phys-ege.sdamgia.ru/problem?id=9306
\BankTask{sdamgia-9306}{
В результате перехода спутника Земли с одной круговой орбиты на другую скорость его движения уменьшается. Как изменяются при этом центростремительное ускорение спутника и период его обращения вокруг Земли?\par1. Увеличивается. 2. Уменьшается. 3. Не изменяется.\par\TaskChangeTable{Центростремительное ускорение спутника}{Период обращения спутника вокруг Земли}
}{
\smallskip\noindent\textbf{Решение}\par
При движении по круговой орбите радиус скорость спутника, его центростремительное ускорение и радиус орбиты связаны соотношением $a= \frac{v ^2}{r} .$ В поле тяжести планеты ускорение равно $a= \frac{GM}{r ^2} ,$ где $M$ — масса планеты. Таким образом, скорость и ускорение связаны соотношением $a=\frac{v^4}{GM}.$ При уменьшении скорости ускорение уменьшается.

Период обращения вокруг Земли равен $T= \frac{2 \pi  r}{v} = \frac{2 \pi  GM}{v ^3} .$ При уменьшении скорости период увеличивается.

 

Ответ: $21$.
}

% Source sdamgia-23230: https://phys-ege.sdamgia.ru/problem?id=23230
\BankTask{sdamgia-23230}{
Ускорение свободного падения на поверхности некоторой сферической однородной планеты равно $g$, а первая космическая скорость для этой планеты равна $v _1.$ Установите соответствие между физическими величинами и формулами, при помощи которых их можно вычислить. 

\TaskMatchingColumns{Физическая величина}{А) радиус планеты\par 
Б) модуль центростремительного ускорения \par 
спутника, летающего вокруг планеты на высоте $h$ \par 
над ее поверхностью по круговой орбите}{Формула}{$1$) $\frac{gv_1^4}{(v_1^2+gh)^2}$\par 
$2$) $\frac{v _1 ^2}{2g}$\par 
$3$) $\frac{v _1 ^2}{g}$\par 
$4$) $\frac{g v _1 ^2}{v _1 ^2 + gh}$}

\begin{center}\small\renewcommand{\arraystretch}{1.25}\begin{tabular}{|p{0.450\linewidth}|p{0.450\linewidth}|}\hline
А & Б\\\hline
 & \\\hline
\end{tabular}\end{center}
}{
\smallskip\noindent\textbf{Решение}\par
А) Из формулы первой космической скорости $v_1=\sqrt{gR}$ найдем радиус планеты: $R= \frac{v _1 ^2}{g} .$

Б) Из равенства сил $F = ma_{\text{цс}} = \frac{m v ^2}{R + h}$ и $F_{\text{т}}=G \frac{Mm}{( R + h ) ^2} ,$ находим, что: 

\[v ^2 = \frac{GM}{R + h} ,\]

откуда:

\[a_{\text{цс}}=\frac{F_{\text{т}}}{m}=\frac{GM}{(R+h)^2}=\frac{v_1^4}{g(R+h)^2}=\frac{v_1^4g}{(v_1^2+gh)^2}.\]

Ответ: $31$.
}

% Source sdamgia-31880: https://phys-ege.sdamgia.ru/problem?id=31880
\BankTask{sdamgia-31880}{
Спутник массой $m$ обращается вокруг Земли по круговой орбите, высота которой над поверхностью планеты равна $h$. Радиус Земли равен $R$. Установите соответствие между физическими величинами и формулами, выражающими их в рассматриваемой задаче ($g$ — ускорение свободного падения на поверхности Земли).

\TaskMatchingColumns{Физическая величина}{А) модуль скорости спутника\par 
Б) период обращения спутника}{Формулы}{$1$) $R\sqrt{\frac{g}{R+h}};$\par 
$2$) $\frac{R\sqrt g}{(R+h)^{3/2}};$\par 
$3$) $\frac{2\pi(R+h)^{3/2}}{R\sqrt g};$\par 
$4$) $\frac{mgR ^2}{2 ( R + r )} .$}

\begin{center}\small\renewcommand{\arraystretch}{1.25}\begin{tabular}{|p{0.450\linewidth}|p{0.450\linewidth}|}\hline
$A$ & Б\\\hline
 & \\\hline
\end{tabular}\end{center}
}{
\smallskip\noindent\textbf{Решение}\par
А) При движении спутника вокруг Земли первая космическая скорость равна $v=\sqrt{\frac{GM}{R+h}},$ при этом ускорение свободного падения $g= \frac{GM}{R ^2} .$ Объединяя формулы, получаем: 

\[(1)\quad v=R\sqrt{\frac{g}{R+h}}.\]

Б) При движении по окружности период обращения $T= \frac{2 \pi  ( R + h )}{v} .$ Объединяя формулы, получаем: 

\[(3)\quad T=\frac{2\pi(R+h)^{3/2}}{R\sqrt g}.\]

Ответ: $13$.
}

% Source sdamgia-19724: https://phys-ege.sdamgia.ru/problem?id=19724
\BankTask{sdamgia-19724}{
Однородная сферическая планета радиусом $R$ имеет плотность $\rho .$ Установите соответствие между физическими величинами и формулами, по которым эти величины можно рассчитать. 

\TaskMatchingColumns{Физическая величина}{А) ускорение свободного падения на поверхности планеты\par 
Б) первая космическая скорость для этой планеты}{Формула}{$1$) $2R\sqrt{\frac{\pi G\rho}{3}}$\par 
$2$) $\frac23R\sqrt{\pi G\rho}$\par 
$3$) $4 \pi  G\rho R$\par 
$4$) $\frac{4}{3} \pi  G\rho R$}

\begin{center}\small\renewcommand{\arraystretch}{1.25}\begin{tabular}{|p{0.450\linewidth}|p{0.450\linewidth}|}\hline
А & Б\\\hline
 & \\\hline
\end{tabular}\end{center}
}{
\smallskip\noindent\textbf{Решение}\par
А) Применим формулы ускорения свободного падения $g = \frac{GM}{R ^2} ,$ массы тела $M = \rho V,$ объема шара $V = \frac{4}{3} \pi  R ^3 .$ В результате получаем общую формулу:

\[g = \frac{4 \pi  G\rho R ^3}{3R ^2} = \frac{4 \pi  G\rho R}{3} .\]

Б) Применим формулу первой космической скорости $v=\sqrt{\frac{GM}{R}},$ и формулы массы и объема. Получим в результате: 

\[v=\sqrt{\frac{4\pi G\rho R^3}{3R}}=2R\sqrt{\frac{\pi G\rho}{3}}.\]

Ответ: $41$.
}

% Source sdamgia-735: https://phys-ege.sdamgia.ru/problem?id=735
\BankTask{sdamgia-735}{
Средняя плотность планеты Плюк равна средней плотности Земли, а радиус Плюка в два раза больше радиуса Земли. Во сколько раз первая космическая скорость для Плюка больше, чем для Земли?
}{
\smallskip\noindent\textbf{Решение}\par
Первая космическая скорость — это скорость, которую необходимо придать объекту, чтобы вывести его на круговую орбиту с радиусом, равным радиусу планеты. Для Земли и Плюка она равна $v_{\text{З}}=\sqrt{\frac{GM_{\text{З}}}{R_{\text{З}}}}$ и $v_{\text{П}}=\sqrt{\frac{GM_{\text{П}}}{R_{\text{П}}}}$ соответственно. Следовательно:

\[\frac{v_{\text{П}}}{v_{\text{З}}}=\sqrt{\frac{M_{\text{П}}R_{\text{З}}}{R_{\text{П}}M_{\text{З}}}}=\sqrt{\frac{\frac43\pi\rho_{\text{П}}R_{\text{П}}^3R_{\text{З}}}{R_{\text{П}}\frac43\pi\rho_{\text{З}}R_{\text{З}}^3}}=\sqrt{\frac{\rho_{\text{П}}R_{\text{П}}^2}{\rho_{\text{З}}R_{\text{З}}^2}}=\sqrt{\frac{(2R_{\text{З}})^2}{R_{\text{З}}^2}}=2.\]

Таким образом, первая космическая скорость на Плюке в два раза больше, чем на Земле.

Ответ: $2$.
}

% Source sdamgia-4439: https://phys-ege.sdamgia.ru/problem?id=4439
\BankTask{sdamgia-4439}{
Известно, что один оборот вокруг своей оси Венера совершает примерно за $243$ земных суток, а масса Венеры составляет $0{,}82$ от массы Земли. На орбиту какого радиуса надо вывести спутник Венеры, чтобы он все время «висел» над одной и той же точкой поверхности? Известно, что спутники Земли, «висящие» над одной и той же точкой поверхности, летают по орбите радиусом $R_{\text{з}}\approx42000\,\text{км}$.
}{
\smallskip\noindent\textbf{Решение}\par
При движении спутника по круговой орбите радиусом $R$ вокруг планеты центростремительное ускорение обеспечивается силой гравитационного притяжения спутника к планете, поэтому второй закон Ньютона записывается в виде $F_{\text{т}} = ma$ или 

\[m\omega ^2 R= \frac{GmM}{R ^2} ,\]

 где $m$ и $M$ — массы спутника и планеты, $G$ — гравитационная постоянная, $\omega = \frac{2 \pi }{T}$ — угловая скорость вращения спутника вокруг планеты, $T$ — период обращения спутника. Подставляя, получаем величину радиуса орбиты:

\[R=\sqrt[3]{\frac{GM}{\omega^2}}=\sqrt[3]{\frac{GMT^2}{4\pi^2}}.\]

Из полученной формулы найдем отношение радиусов орбит спутников Венеры и Земли: 

\[\frac{R_{\text{В}}}{R_{\text{З}}}=\sqrt[3]{\frac{GM_{\text{В}}T_{\text{В}}^2}{4\pi^2}}:\sqrt[3]{\frac{GM_{\text{З}}T_{\text{З}}^2}{4\pi^2}}=\sqrt[3]{\frac{M_{\text{В}}}{M_{\text{З}}}\cdot\frac{T_{\text{В}}^2}{T_{\text{З}}^2}}.\]

Период обращения спутника вокруг Земли $T_{\text{З}} = 1$ сутки, поэтому, подставляя данные из условия, находим:

\[R_{\text{В}}=\sqrt[3]{0{,}82\cdot243^2}R_{\text{З}}\approx1531000\,\text{км}.\]

Ответ: $R_{\text{В}}=R_{\text{З}}(0{,}82\cdot243^2)^{1/3}\approx1531000\,\text{км}.$

Примечание Д. Д. Гущина.

Внимательный читатель заметит, что найденный радиус орбиты превышает радиус сферы Хилла, равный для Венеры одному миллиону километров. Иными словами, притяжения Венеры не хватает, чтобы удерживать настолько удаленный от нее спутник. Поэтому со временем он покинет орбиту Венеры и станет независимо от Венеры вращаться вокруг Солнца. Чтобы этого не произошло, необходимо (пока позволяют запасы топлива) корректировать орбиту спутника двигателями.

Пытливый читатель может заинтересоваться вопросом о том, каков радиус сферы Хилла для Земли. Он равен $1{,}47\,\text{млн\,км}$. С удовольствием отметим, что Луна удалена от Земли на $0{,}384\,\text{млн\,км}$, поэтому Луна не покинет орбиту Земли. Подробности об этом можно прочитать в англоязычной Википедии.
}

% Source sdamgia-31899: https://phys-ege.sdamgia.ru/problem?id=31899
\BankTask{sdamgia-31899}{
$22$ октября 2022 года с космодрома «Восточный» ракетой из семейства «Союз» были запущены три спутника связи «Гонец-М» и первый спутник «Скиф-Д» из проекта «Сфера» по обеспечению доступа к широкополосному интернету на всей территории России. Этот спутник вышел на круговую орбиту высотой $h=8170\,\text{км}$ над поверхностью Земли с полярным наклоном $90^\circ$ (плоскость такой орбиты проходит через оба географических полюса Земли), и через некоторое время пролетел с юга на север над центром Москвы. Через сколько оборотов вокруг Земли он снова окажется над Москвой, двигаясь в том же направлении с юга на север? Радиус Земли $R_{\text{з}}=6370\,\text{км}$.
}{
\smallskip\noindent\textbf{Решение}\par
$1$. Плоскость орбиты спутника сохраняет свое положение в пространстве, а Земля вертится вокруг оси, и точки на ее поверхности пересекают эту плоскость по два раза за период обращения спутника, причем нас интересует только одно пересечение, когда спутник пролетает над Москвой в направлении с юга на север.

$2$. Чтобы спутник снова пролетел над Москвой в том же направлении, очевидно, необходимо выполнение соотношения $mT_{\text{з}} = nT_{\text{сп}},$ где $m$ и $n$ — целые числа, а $T_{\text{з}}=24\,\text{часа}=1440\,\text{мин}$ и $T_{\text{сп}}$ — периоды обращения Земли вокруг своей оси и спутника по круговой полярной орбите вокруг Земли.

$3$. Для нахождения $T_{\text{сп}}= \frac{2 \pi  r}{v _1} ,$ где $r=R_{\text{з}}+h$, а $v _1$— скорость спутника массой $m_{\text{с}}$, запишем для него второй закон Ньютона с учетом выражений для силы тяжести Земли $F= \frac{GMm_{\text{с}}}{r ^2}$ и центростремительного ускорения $a= \frac{v _1 ^2}{r}$:

\[\frac{GMm_{\text{с}}}{r ^2} = \frac{m_c v _1 ^2}{r} ,\]

 откуда $v_1=\sqrt{\frac{GM}{r}},$ и поскольку $g_0= \frac{GM}{R ^2 _{\text{з}}} ,$ то:

\[T_{\text{сп}}=2\pi r\sqrt{\frac{r}{GM}}=\frac{2\pi}{R_{\text{з}}}\sqrt{\frac{(R_{\text{з}}+h)^3}{g_0}}.\]

 Здесь $g_{0}$ и $g_1$ — ускорения свободного падения на поверхности Земли и на высоте $h$ над ней.

$4$. Подставляя численные данные из условия и таблиц, получаем:

\[T_{\text{сп}}=\frac{2\cdot3{,}14}{6370000}\sqrt{\frac{(6370000+8170000)^3}{10}}\approx17285\,\text{с}\approx288{,}08\,\text{мин}.\]

$5$. Таким образом:

\[\frac{n}{m} = \frac{T_{\text{з}}}{T_{\text{сп}}} \approx \frac{1440}{288{,}08} \approx 4{,}999\approx 5,\]

 то есть за одни сутки спутник совершит $5$ оборотов вокруг Земли и снова окажется над Москвой.

Ответ: $5$ оборотов.
}


% Source F64C40: https://ege.fipi.ru/bank/questions.php?search=1&proj=BA1F39653304A5B041B656915DC36B38&qid=F64C40
\BankTask{F64C40}{
Искусственный спутник обращается вокруг Земли по вытянутой эллиптической орбите. В некоторый момент времени спутник проходит точку минимального удаления от Земли.
\begin{enumerate}
\item Кинетическая энергия спутника при прохождении этого положения максимальна.
\item Сила притяжения спутника к Земле в этом положении максимальна.
\item Полная механическая энергия спутника в этом положении достигает максимума.
\item Скорость спутника при прохождении этого положения минимальна.
\item Ускорение спутника при прохождении этого положения равно $0$.
\end{enumerate}
}{
\smallskip\noindent\textbf{Решение от GPT}\par
В перигее $r$ минимален, поэтому $F=GMm/r^2$ максимальна. Полная энергия постоянна; потенциальная $-GMm/r$ минимальна, следовательно, кинетическая энергия и скорость максимальны. Ответ: $12$.
}

% Source ACFD5A: https://ege.fipi.ru/bank/questions.php?search=1&proj=BA1F39653304A5B041B656915DC36B38&qid=ACFD5A
\BankTask{ACFD5A}{
Искусственный спутник обращается вокруг Земли по вытянутой эллиптической орбите. В некоторый момент времени спутник проходит положение максимального удаления от Земли.
\begin{enumerate}
\item Скорость спутника при прохождении этого положения максимальна.
\item Сила притяжения спутника к Земле при прохождении этого положения максимальна.
\item Полная механическая энергия спутника при прохождении этого положения максимальна.
\item Импульс спутника при прохождении этого положения минимален.
\item Ускорение спутника при прохождении этого положения минимально.
\end{enumerate}
}{
\smallskip\noindent\textbf{Решение от GPT}\par
В апогее $r$ максимален; поэтому $F=GMm/r^2$ и $a=GM/r^2$ минимальны. Полная механическая энергия постоянна, потенциальная максимальна, значит скорость и импульс минимальны. Ответ: $45$.
}

% Source F4D3FF: https://ege.fipi.ru/bank/questions.php?search=1&proj=BA1F39653304A5B041B656915DC36B38&qid=F4D3FF
\BankTask{F4D3FF}{
Космический аппарат, обращающийся вокруг Луны по круговой орбите, перешёл на другую круговую орбиту бóльшего радиуса. Как изменились в результате этого перехода центростремительное ускорение, с которым аппарат движется по орбите, и его период обращения вокруг Луны?
\begin{enumerate}
\item увеличилась
\item уменьшилась
\item не изменилась
\end{enumerate}
\TaskChangeTable{Центростремительное ускорение}{Период обращения аппарата вокруг Луны}
}{
\smallskip\noindent\textbf{Решение от GPT}\par
$a=GM/r^2$ уменьшается с ростом радиуса; $T=2\pi\sqrt{r^3/(GM)}$ увеличивается. Ответ: $21$.
}

\end{document}
